\documentclass[11pt]{article}

\usepackage[margin=1in]{geometry}
\usepackage{amsmath,amssymb}
\usepackage{booktabs}
\usepackage{array}
\usepackage{pifont}
\usepackage[colorlinks=true,citecolor=blue,linkcolor=blue,urlcolor=blue]{hyperref}
\usepackage[round,authoryear]{natbib}
\usepackage{times}

\newcommand{\eps}{\varepsilon}

\title{Related Work: Positioning FlowSurv-AFT in the Landscape of\\
Survival Analysis Methodology}
\author{}
\date{}

\begin{document}
\maketitle

\section{Introduction to the review}

Time-to-event (survival) analysis is distinguished from ordinary regression
by \emph{censoring}: for a subset of subjects the event of interest --
death, relapse, mechanical failure, default -- has not occurred by the end
of follow-up, so only a lower bound on the true event time is observed. A
survival model must therefore do two things simultaneously: incorporate
censored observations correctly into estimation, and represent the shape of
the conditional event-time distribution, whose hazard function may be
monotone increasing (wear-out), monotone decreasing (early-life failure),
bathtub-shaped (high early risk, a stable middle period, rising late risk --
canonical in reliability engineering and in oncology follow-up), multimodal
(heterogeneous sub-populations pooled under one label), or non-proportional
across covariate strata (treatments whose relative benefit reverses in
time). No published model, to our knowledge, combines an arbitrary
learned hazard shape, exact closed-form inference in both directions, and
the interpretability of the classical accelerated failure time (AFT)
framework, estimated by the full right-censored likelihood. This section
develops that claim in detail, tracing four methodological lineages --
classical, machine-learning, discriminative/discrete-time deep, and
generative/flow-based -- to the point where each is shown to surrender
at least one of these four properties, and situates the present
contribution, \emph{FlowSurv-AFT}, in the resulting gap.

The review is organized as follows. Section~\ref{sec:classical} covers the
classical parametric and semiparametric baselines that remain the
default choice in applied practice. Section~\ref{sec:ml} covers
ensemble/machine-learning survival models and the recent large-scale
neutral-comparison evidence that reframes what a "win" should mean in
this literature. Section~\ref{sec:dl} reviews deep survival models that
retain a proportional-hazards or discrete-time structure.
Section~\ref{sec:deepaft} is a focused critical analysis of the deep AFT
literature, which we argue has consistently traded away the exact
likelihood in exchange for a flexible covariate map.
Section~\ref{sec:generative} covers generative and flow-based survival
models, with particular emphasis on the two works closest to the present
contribution -- \citet{ausset2021flows} and the concurrent
\citet{licai2026} -- which we differentiate in detail.
Section~\ref{sec:evaluation} reviews the evaluation methodology literature
that constrains which claims this line of work is entitled to make.
Section~\ref{sec:synthesis} synthesizes the review into a four-camp
taxonomy and states the resulting research gap formally.

\section{Classical survival models}
\label{sec:classical}

The Cox proportional-hazards model \citep{cox1972} remains survival
analysis's default tool. It is semiparametric: the partial likelihood
estimates the log-hazard-ratio coefficients $\beta$ without specifying the
baseline hazard $h_0(t)$, which is recovered only as a non-smooth
Breslow-type step function if a survival curve is required at all. Two
consequences follow directly from this construction. First, Cox regression
assumes proportional hazards -- $h(t\mid x) = h_0(t)\exp(\beta^\top x)$ --
so it cannot represent hazard ratios that vary in time (crossing survival
curves), without extension (time-varying coefficients, stratification).
Second, because the baseline hazard is nonparametric, the model yields no
smooth conditional density $f(t\mid x)$ and no closed-form quantile
function; both must be obtained, if at all, by post-hoc smoothing of the
step-function estimate.

The accelerated failure time (AFT) model \citep{wei1992aft} takes the
complementary route: it specifies a full parametric law for $\log T$,
\begin{equation}
  \log T = \beta^\top x + \sigma \eps, \qquad \eps \sim F_0,
  \label{eq:aft}
\end{equation}
with $F_0$ a fixed standard distribution -- extreme-value (giving a Weibull
$T$), logistic (log-logistic $T$), or normal (log-normal $T$). This
buys two things a Cox model does not provide: a closed-form density,
survival, hazard and quantile function in one pass, and the
\emph{acceleration factor} interpretation, $\mathrm{TR} =
\exp\{\beta^\top(x_a - x_b)\}$, the multiplicative change in the entire
event-time distribution (every quantile simultaneously) between two
covariate profiles -- an effect measure that is arguably more intuitive to
non-statistician stakeholders than a hazard ratio, and the natural
communication tool in reliability engineering, where "time-to-failure is
multiplied by $\mathrm{TR}$" maps directly onto warranty and maintenance
planning. The AFT's cost is the parametric assumption on $F_0$: a
Weibull hazard $h(t) = \lambda\rho(\lambda t)^{\rho-1}$ is monotone in $t$
for every $\rho \ne 1$ and constant only at $\rho=1$; it can never turn
over, so it is structurally incapable of representing a bathtub or a
multimodal hazard, regardless of how much data is available to estimate
$\beta$. The log-normal AFT is more flexible (a hump-shaped hazard) but
is likewise unimodal and cannot represent a bathtub.

\citet{royston2002}'s flexible parametric model relaxes the shape
constraint by modelling the log cumulative hazard as a restricted cubic
spline in $\log t$ rather than a fixed parametric form,
\[
  \log H(t\mid x) = s(\log t;\, \gamma) + \beta^\top x,
\]
with knots placed at quantiles of the observed event log-times. This is
the strongest classical shape-flexible baseline available and the natural
comparator for any deep flexible-hazard model; because the spline is a
free-form nonparametric object rather than a probability-generating
transform of a single latent variable, however, it does not carry the
location-scale AFT interpretation of Eq.~\eqref{eq:aft} unless further
restricted, and its flexibility is capped by a small, hand-chosen number
of degrees of freedom rather than learned end-to-end from covariates.
\citet{opatha2024} provide the classical anchor and direct motivation for
the present work: a careful simulation study mapping Cox-versus-Weibull-AFT
performance across sample size, censoring proportion, censoring type
(Type~I administrative vs.\ Type~III random), and hazard shape, showing that
the AFT dominates specifically at small sample sizes and under monotone
increasing hazards, and that the ranking is sensitive to exactly the design
factors this study extends to deep learning. Their study, however, was
confined to two classical models and to Harrell's concordance, leaving
open where deep, flexible-density models sit on the same map, and under
which modern calibration instruments (Section~\ref{sec:evaluation}) --
the question this work takes up.

\section{Machine-learning survival models and the neutral-comparison evidence}
\label{sec:ml}

Random Survival Forests \citep{ishwaran2008rsf} and gradient-boosted
Cox variants extend Cox regression's covariate flexibility -- nonlinearity,
interaction detection -- without positing a hazard family, at the cost of a
smooth conditional density: RSF's native output is an ensemble-averaged
Nelson--Aalen-type cumulative hazard estimate, again a step function, with
only weak calibration guarantees. The decisive fact reshaping how any
gain from this literature must be argued, however, is empirical rather
than architectural: \citet{burk2024neutral}, in the largest neutral
comparison to date (19 methods across 34 datasets, all methods tuned
under a common protocol so that no method is home-advantaged), find that
\emph{no method -- including every machine-learning and deep-learning
method surveyed -- significantly outperforms plain Cox regression in
discrimination on standard low-dimensional tabular data}. This finding is
the single most important piece of evaluation evidence this review draws
on: it means a flexible deep survival model's case cannot rest on
concordance improvements, which the neutral evidence says are not there to
be found on this class of data. The contribution this literature is
capable of making instead is calibration and distributional fidelity --
whether the \emph{shape} of the predicted individual survival curve,
not merely its rank ordering, matches the truth -- and that reframing
directly motivates the evaluation design in Section~\ref{sec:evaluation}
and the choice, throughout this work, never to lead with a concordance
claim.

\section{Deep learning survival models}
\label{sec:dl}

Deep architectures for survival analysis divide naturally into three
families by what structural assumption they retain from the classical
models above.

\subsection{Cox-family deep models}
DeepSurv \citep{katzman2018deepsurv} is the most direct deep extension of
Cox regression: it replaces the linear predictor $\beta^\top x$ inside the
Cox partial likelihood with a feed-forward network, gaining nonlinearity
and interaction modelling but retaining every structural limitation of the
Cox model itself -- proportional hazards, a Breslow-estimated baseline, no
smooth density. Cox-Time \citep{kvamme2019coxtime} generalizes further by
allowing time--covariate interaction terms, which lets the model escape
strict proportionality, but the estimator still factorizes through a
shared, non-parametrically estimated baseline hazard scaled by a
covariate function, so the output remains a step-function survival curve
fit by partial, not full, likelihood.

\subsection{Discrete-time distribution models}
A second family discretizes time onto a grid and fits a
distribution over grid cells directly, sidestepping the Cox
baseline-hazard machinery entirely. DeepHit \citep{lee2018deephit} places
a softmax output over a fixed set of time bins, giving a genuine
per-subject probability mass function (and natively handling competing
risks); DRSA \citep{ren2019drsa} predicts the same object
autoregressively, one grid step at a time; PC-Hazard
\citep{kvamme2021pchazard} instead models a piecewise-constant hazard on
each grid interval, which admits a proper (if piecewise) likelihood
derivation; SurvTRACE \citep{wang2022survtrace} replaces the covariate
encoder with a transformer but keeps the discrete-hazard output head. All
four are shape-free \emph{within the resolution of the chosen grid}, and
all four inherit the same structural cost: the predicted survival curve is
a step function by construction, so density and hazard exist only as
grid-cell probabilities (not evaluable at an arbitrary $t$), and quantiles
-- including the median survival time used for concordance -- are only
available up to the grid's bin width. Grid placement is consequently not a
cosmetic choice but a bias--variance trade-off baked into the estimator:
a coarse grid discards temporal resolution the data may support; a fine
grid multiplies the number of estimated parameters and can destabilize
training under heavy censoring, where many bins receive few or no events.
We refer to this as the discretization horn throughout this work.

\subsection{Continuous-time flexible models}
A smaller family attempts to retain continuous time without discretizing.
SODEN \citep{tang2022soden} generalizes the DeepSurv/DeepHit family
through a system of ordinary differential equations whose solution
implicitly defines the hazard, buying flexibility at the cost of a
numerically delicate, implicitly-defined density that must be solved for
at every evaluation. SuMo-net \citep{rindt2022sumonet} instead
parameterizes $S(t\mid x)$ directly with a network constrained to be
monotone decreasing in $t$, trained with a proper scoring rule; this gives
smooth, continuous survival curves, but the corresponding density and
hazard functions exist only through numerical differentiation of the
learned monotone map, which is not exact and amplifies estimation noise
precisely in the tails where hazard-shape questions (bathtub, multimodal)
are decided.

\section{Deep accelerated failure time models: a critical synthesis}
\label{sec:deepaft}

The AFT form's modern revival under deep learning is real, but -- we argue
after close reading of each contribution -- structurally incomplete in a
single, consistent way. \emph{deepAFT} \citep{norman2024deepaft} places a
deep network on the location term $\mu(x)$ of Eq.~\eqref{eq:aft} while
leaving the residual distribution unspecified, and estimates the model by
Buckley--James iterative imputation, inverse-probability-of-censoring
weighting (IPCW), or a rank/transformation loss rather than by maximizing
a likelihood. This choice is consequential in exactly the setting this
work is designed to stress: IPCW-based estimators are well known to
become unstable as the censoring proportion grows, because the inverse
weights blow up as the censoring-survival estimate approaches zero, and
Buckley--James imputation similarly degrades as the fraction of
uninformative (heavily censored) observations increases. The resulting
survival curves are Kaplan--Meier-rescaled step functions, not smooth
conditional densities. DeepR-AFT \citep{kim2022deepraft} trains with a
Gehan-type rank loss, which supplies a valid concordance-style ranking
objective but, being rank-based, carries no information about the scale
or shape of the predicted distribution and cannot produce a calibrated
density at all. KAN-AFT \citep{francis2025kanaft} substitutes a
Kolmogorov--Arnold network for the covariate map, trading the deep
network's opacity for (claimed) symbolic interpretability of $\mu(x)$,
but inherits the same likelihood-free estimation as deepAFT. The most
recent likelihood-free entrant identified in our updated sweep,
\emph{GRAFT} \citep{graft2026}, follows the identical pattern -- a
gated residual AFT architecture trained with a concordance-aligned rank
loss and post-hoc calibration correction -- and \emph{DART}
\citep{dart2023} likewise trains a deep AFT with a Gehan-type rank
objective; both therefore occupy exactly the same cell of the taxonomy as
deepAFT/DeepR-AFT/KAN-AFT and do not alter the gap analysis below.

Two contributions do retain the likelihood, but at the cost of the
residual law's flexibility. Neural Conditional Event Time models
\citep{engelhard2020neuralcet} -- Neural-CET -- place neural networks on
\emph{both} the location $\mu(x)$ and the scale $\sigma(x)$ and estimate
by the exact censored likelihood, which is structurally the closest prior
work to the present model; the residual law is nonetheless fixed to a
standard Gaussian, so the model is, in the terminology of
Eq.~\eqref{eq:aft}, a log-normal AFT with a neural covariate map -- shape
constrained to the log-normal's single hump, unable to represent a
bathtub or a multimodal hazard by construction, whatever the capacity of
the covariate network. This is precisely the ablation this study
isolates as \emph{FlowSurv-Gauss}: freeze the flow to the identity and fix
the base law to Gaussian, and Neural-CET's specification is recovered
exactly, which lets the empirical study attribute any performance gap
directly to the residual flow rather than to encoder capacity or training
protocol differences. DKAFT \citep{wu2021dkaft} similarly retains the
likelihood but represents the residual through a Gaussian-process output
layer -- again log-normal in effect. \citet{ranganath2016deepsurv}'s deep
exponential family parameterizes Weibull event times through a deep
generative model, retaining the Weibull's monotone-hazard restriction
even as the covariate mapping deepens.

Taken together, the deep AFT literature exhibits a consistent trade-off:
\emph{every published deep AFT model that we are aware of either keeps
the exact censored likelihood but fixes the residual family to a
single-mode distribution (Gaussian in Neural-CET/DKAFT, Weibull in
Ranganath et al.), or frees the covariate map and the estimation
objective but abandons the likelihood entirely (deepAFT, DeepR-AFT,
KAN-AFT, GRAFT, DART).} No prior work combines a free, learned residual
law with exact full-likelihood estimation inside the AFT form; this is
the specific gap FlowSurv-AFT is designed to close.

\section{Generative and flow-based survival models}
\label{sec:generative}

\subsection{Mixture-based densities}
Deep Survival Machines \citep{nagpal2022dsm} -- DSM -- parameterizes a
finite mixture of $K$ Weibull or log-normal components, with an MLP
mapping covariates to per-component mixture weights and parameters, and
is trained by the exact censored likelihood; it is consequently the
standard fully-parametric, exact-likelihood baseline against which this
work benchmarks throughout. Deep Cox Mixtures \citep{nagpal2021coxmixtures}
and the earlier DeepWeiSurv/DPWTE line
\citep{bennis2020deepweisurv,bennis2021dpwte} are architectural
predecessors and variants of the same idea. Survival MDN
\citep{han2022survivalmdn} instead places a Gaussian mixture on a
transformed time axis, deriving the survival function by an explicit
change of variables. The shared limitation, acknowledged directly by the
Survival MDN authors themselves, is that every one of these constructions
builds a multimodal or bathtub-shaped hazard, if at all, out of a
\emph{finite sum of unimodal, monotone-or-hump-shaped primitives}: with
few mixture components the model cannot represent a genuine bathtub;
with many components the number of mixture weights to fit under
censoring grows, $K$ becomes a fragile hyperparameter, and the
resulting hazard, while technically flexible in the infinite-$K$ limit, is
not a \emph{learned} shape in the same sense a normalizing flow's
residual law is -- it is a discrete architectural choice made in advance
of seeing how much flexibility the data actually needs.

\subsection{Adversarial generators}
\citet{chapfuwa2018adversarial} and the subsequent survival-function
matching formulation \citep{chapfuwa2021sfm}, together with related work
by Zhou et al., generate event times through an adversarial (GAN-style)
mechanism, which in principle imposes no shape restriction on the
generated distribution at all. The cost is that adversarial training
supplies no tractable likelihood and no closed-form survival function --
both must be estimated empirically from generated samples -- and GAN
training is well documented to be unstable relative to direct
likelihood maximization, a particular liability under the class imbalance
that heavy censoring induces.

\subsection{Normalizing flows for survival: the direct lineage}
\label{sec:flows}
The line of work most directly antecedent to FlowSurv-AFT applies
normalizing flows -- invertible, density-tracking transformations of a
simple base distribution -- to survival times. \citet{miscouridou2018} is
the earliest instance we are aware of: a discrete invertible transform of
a Weibull base distribution, but embedded inside a latent-variable model
trained by black-box variational inference, so the fitted objective is an
evidence lower bound (ELBO), not the exact right-censored likelihood that
the present work and its comparators require.

\citet{ausset2021flows} (elaborated in the accompanying doctoral thesis,
\citealp{ausset2021thesis}) is the closest prior work and the paper this
study differentiates most carefully against. Ausset et al.\ condition a
\emph{continuous} normalizing flow -- a neural ordinary differential
equation in the FFJORD family -- on covariates and on log survival time,
and train it with the exact right-censored likelihood, remarking in
passing that the construction "can be seen as a generalization of the AFT
model." Having examined the thesis in detail, we verify the following.
First, every evaluation of the density or the survival function requires
solving an initial-value ODE numerically; there is no closed form for
$f(t\mid x)$, $S(t\mid x)$, a quantile, or a sample, and the discretization
error of the ODE solver becomes a further, uncontrolled source of
approximation error layered on top of the statistical estimation error.
Second, the thesis reviews discrete coupling and autoregressive
normalizing flows as alternatives and explicitly sets them aside in favour
of the continuous formulation, so the analytic-inverse route this work
takes was a known, considered, and declined alternative rather than an
oversight. Third, the empirical evaluation is comparatively narrow: one
multimodal synthetic scenario at a single, fixed 80\% censoring level, with
Harrell's C as the only reported metric -- no hazard-shape recovery
error against ground truth, no calibration analysis (D-calibration, ICI,
or any proper scoring rule), and no controlled sweep over censoring
proportion or censoring mechanism. Fourth, and most directly relevant to
the interpretability claim of the present work, the "generalization of
the AFT" remark is a one-sentence framing aside: the fitted model has no
explicit $\mu(x)/\sigma(x)$ decomposition and reports no time-ratio
quantity, so no acceleration-factor interpretation is extracted or
validated anywhere in the thesis. FlowSurv-AFT differs from this work
along exactly these four axes: it is bidirectionally exact (no ODE
solver, closed-form in both directions), it carries an explicit,
validated AFT decomposition, and it is evaluated by the hazard-shape and
calibration protocol Ausset et al.\ do not report.

The second closest work is concurrent and, at the time of writing,
unpublished: \citet{licai2026}, under review for STAI-X 2026. Li \& Cai
build \emph{discrete monotone tanh-sum flows} directly on log event time,
trained with an exact forward (density and survival) right-censored
likelihood, augmented by a Soft--Nelson--Aalen Wasserstein regularizer and
a shared-latent extension addressing informative censoring. Having
reviewed the submitted manuscript under the project's pre-registered
novelty-sweep protocol, we verify three properties that jointly
distinguish it from FlowSurv-AFT. First, the tanh-sum transform has no
closed-form inverse: every quantile evaluation or sample draw requires an
iterative 40-step numerical bisection procedure, so while the
\emph{forward} (data-to-likelihood) direction is exact and closed-form,
the \emph{reverse} direction that bidirectional exactness requires is not
-- this is the concurrent work's structural analogue of Ausset et al.'s
ODE-solver dependency, arrived at from the opposite (discrete rather than
continuous) direction. Second, the transform's knot parameters are
themselves general functions of the covariates $x$ rather than entering
through a location-scale form as in Eq.~\eqref{eq:aft}, so the map is not
of AFT type and admits no acceleration-factor interpretation. Third, the
reported simulation study uses only standard parametric distribution
families for the data-generating process, with no bathtub, multimodal, or
crossing-hazard scenario, no calibration metric, and no controlled
censoring-proportion sweep. The distribution-matching regularizer and the
informative-censoring extension are genuine contributions in their own
right and are cited here as complementary, not competing, technical
devices; the empty methodological cell this study occupies -- exact
bidirectional tractability \emph{and} an explicit AFT decomposition
\emph{and} the hazard-shape/calibration-across-censoring study -- remains
open in this concurrent manuscript as reviewed.

\subsection{Transformation models: the statistical-equivalence view}
A parallel statistical literature arrives at essentially the same object
-- a flexible, invertible, monotone transform of a latent variable,
trained by a censored likelihood -- from distributional regression rather
than from the normalizing-flow literature. Deep conditional
transformation models \citep{baumann2021transformation} and DRIFT
\citep{kolb2024drift} establish the formal equivalence between a
one-dimensional conditional normalizing flow and a flexible
distributional-regression model fit by a censored likelihood: these are,
statistically, the same estimator viewed through a different vocabulary,
and we cite them as the relevant related work from that vocabulary rather
than as a distinct methodological camp. A further recent entrant in this
family identified in the updated sweep, \citet{yuan2025transformation},
develops an interval-censored partially-linear transformation model fit
by monotone-spline sieve maximum likelihood combined with a deep neural
covariate map; it sits in the same transformation-model camp as
DRIFT/deep conditional transformation models -- flexible and
likelihood-based, but not cast in AFT location-scale form and not
evaluated for hazard-shape recovery across a censoring-proportion sweep.

\subsection{Dependent censoring: an explicit scope boundary}
\citet{gharari2023copula} address the harder and distinct problem of
\emph{dependent} (informative) censoring using deep copula models with
identifiability guarantees. Every model discussed in this review,
including FlowSurv-AFT, assumes censoring is independent of the event
time given covariates; dependent censoring is explicitly out of scope
here and is noted as a limitation (Section~\ref{sec:synthesis}), with
Gharari et al.\ as the relevant reference point for extending this line
of work in that direction.

\subsection{Adjacent 2025--2026 landscape}
A final literature sweep, frozen with the project's pre-registration on
2026-08-04, searched arXiv, OpenReview, and Semantic Scholar for any
post-2024 work at the intersection of normalizing flows, splines, and
survival analysis. Beyond \citet{ausset2021flows} and the concurrent
\citet{licai2026}, no spline-flow survival model and no new
exact-likelihood flow competitor was found. The nearest adjacent items --
included here for completeness and to substantiate that the sweep was
thorough rather than narrowly targeted -- are: \citet{lticl2026}, a
conditional normalizing-flow head for right-censored \emph{lead-time
forecasting} in an industrial supply-chain, in-context-learning setting,
sharing the "conditional flow under right censoring" idea but with no AFT
decomposition, no spline flow, and evaluated by CRPS in a non-clinical,
non-reliability domain; \citet{yin2025multimodal}, which uses a
normalizing flow for cross-modal representation alignment in multimodal
whole-slide-image survival prediction rather than as the event-time
model itself; \citet{graft2026} and \citet{dart2023}, already placed in
the likelihood-free deep AFT camp above (Section~\ref{sec:deepaft});
\citet{yuan2025transformation}, placed in the transformation-model camp
above; \citet{cennsurv2025}, a neural network for exposure-lag survival
modelling -- a distinct problem formulation; \citet{gregorio2023}, a
Royston--Parmar-style spline neural network for heterogeneous
treatment-effect estimation rather than a general-purpose survival model;
and \citet{eguchi2026}, a monograph on flow matching for statistical
inference with a survival-analysis chapter, cited here as context. None
of these alters the gap analysis of Section~\ref{sec:synthesis}.

\section{Evaluation methodology: instruments and pitfalls}
\label{sec:evaluation}

The metrics used to evaluate a survival model shape, in practice, which
methodological trade-offs look favourable, so the evaluation literature is
itself part of the argument this study makes. \textbf{Discrimination.}
Harrell's concordance index is ubiquitous but is known to be
censoring-dependent (it is biased under heavy censoring because
concordance can only be assessed on pairs where the earlier time is
observed); Uno's C \citep{uno2011} corrects this with inverse-probability
weighting, at the cost of requiring a consistent estimate of the
censoring distribution. \citet{sonabend2022chacking} names and formalizes
the risk this study takes care to avoid -- \emph{C-hacking}, reporting
whichever of several concordance variants is most favourable to the
proposed method -- and recommends pre-specifying a single variant; this
study reports Uno's C only, fixed by pre-registration before any
experiment was run, and explicitly declines to claim concordance
superiority as a headline result, consistent with the neutral-comparison
evidence of Section~\ref{sec:ml}. \textbf{Overall accuracy.} The
Integrated Brier Score \citep{graf1999ibs}, with the IPCW correction of
\citet{gerds2006}, summarizes squared-error calibration and discrimination
jointly over a time horizon. \textbf{Calibration.} D-calibration and
1-calibration \citep{haider2020} test whether the predicted survival
probabilities are distributionally correct (a 10\% predicted survival
probability should be observed roughly 10\% of the time, marginally); the
integrated calibration index \citep{austin2020ici} quantifies calibration
error at a fixed horizon via a smoothed calibration curve; X-CAL
\citep{goldstein2020xcal} proposes calibration as an explicit,
differentiable training objective rather than only a post-hoc diagnostic.
\textbf{Position papers.} \citet{lillelund2025} argue directly that the
metric reported must be matched to the claim being made and that the
censoring assumption underlying each metric should be stated explicitly
-- a discipline this study's pre-registration is designed to enforce
mechanically. \citet{pearce2022benchmark} formalize the simulation-rigor
protocol this study's design follows: synthetic data with synthetic,
controllable censoring for ground-truth recovery claims, combined with
real data under real (uncontrolled) censoring for external validity, over
repeated splits rather than a single train/test division. Finally,
\citet{burk2024neutral} (already introduced in Section~\ref{sec:ml}) is
the empirical anchor that converts all of the above from methodological
best practice into a substantive constraint on what this study is
entitled to claim: because large-scale neutral evidence finds no
discrimination advantage for flexible or deep methods on standard tabular
data, the calibration and hazard-shape-recovery instruments reviewed in
this section -- not concordance -- are the evidentiary standard this
study is held to.

Two further methodological references are technical rather than
survival-specific but underpin the model directly. \citet{durkan2019nsf}
introduce the rational-quadratic spline coupling transform whose
monotone, per-bin construction gives FlowSurv-AFT's residual flow its
bidirectional exactness: forward and inverse maps, and their Jacobian
determinants, are all available in closed form within each spline bin,
with identity behaviour imposed outside a bounded domain. \citet
{wiegrebe2024review} and \citet{drysdale2022survset} are general
reference works -- a broad review of deep learning for survival analysis
and an open repository of benchmark time-to-event datasets, respectively
-- cited for context rather than as direct methodological comparators.

\section{Synthesis: the research gap}
\label{sec:synthesis}

The literature reviewed above separates cleanly into four camps, each
camp surrendering exactly one of the four properties a model must have to
combine flexibility, tractability, interpretability, and correct
likelihood-based estimation under censoring. Table~\ref{tab:gap} makes
this explicit.

\begin{table}[h!]
\centering
\footnotesize
\setlength{\tabcolsep}{4pt}
\begin{tabular}{p{3.4cm}p{1.9cm}p{2.5cm}p{1.5cm}p{1.9cm}p{2.7cm}}
\toprule
\textbf{Camp} & \textbf{Smooth, shape-free density} & \textbf{Exact $f,S$, quantiles, sampling (no solver/grid)} & \textbf{AFT interp.} & \textbf{Exact censored likelihood} & \textbf{Shape/calibration study across censoring} \\
\midrule
Classical AFT / Royston--Parmar
  & \ding{55}\,(restricted shapes)
  & \checkmark
  & \checkmark
  & \checkmark
  & partial (\citealp{opatha2024} -- classical models only) \\
Discrete-time DL (DeepHit family)
  & \ding{55}\,(grid steps)
  & \ding{55}
  & \ding{55}
  & \checkmark\,(discrete)
  & \ding{55} \\
Deep AFT (deepAFT, DeepR-AFT, KAN-AFT, GRAFT, DART)
  & \ding{55}\,(no density)
  & \ding{55}
  & \checkmark
  & \ding{55}
  & \ding{55} \\
Generative (DSM/MDN; Ausset; Li \& Cai)
  & \checkmark\,/\,partial
  & \ding{55}\,(mixture-shape limits; ODE solver; bisection inverse)
  & \ding{55}
  & \checkmark
  & \ding{55} \\
\bottomrule
\end{tabular}
\caption{The four-camp taxonomy of the survival-modelling literature
reviewed in this section. Each camp satisfies at most four of the five
properties; the conjunction of all five (rightmost four columns plus a
learned, shape-free density) is, to our knowledge, unoccupied.}
\label{tab:gap}
\end{table}

\emph{The empty cell is the conjunction.} FlowSurv-AFT is constructed to
occupy it directly, by combining (i) an explicit AFT location-scale
decomposition, $\log T = \mu(x) + \sigma(x)\eps$, with (ii) a conditional
rational-quadratic spline residual flow for $\eps$ -- analytically
invertible in both directions, with closed-form Jacobians, following
\citet{durkan2019nsf} -- trained by (iii) the exact right-censored full
likelihood (density term for events, survival term for censored
observations, with no inverse-probability weighting, no partial
likelihood, and no time discretization), and validated by (iv) the first,
to our knowledge, systematic hazard-shape-recovery and calibration study
conducted across a controlled grid of censoring proportion and censoring
mechanism, extending the classical simulation design of
\citet{opatha2024} into this flexible-density setting. Sections
\ref{sec:deepaft} and~\ref{sec:flows} establish that no prior
contribution combines properties (i)--(iii); this section's synthesis,
together with the frozen 2026-08-04 sweep of Section~\ref{sec:generative}
(``Adjacent 2025--2026 landscape''), establishes that this remains true as
of the most recent literature check preceding the pre-registration of
this study's hypotheses and confirmatory contrasts.

\bibliographystyle{apalike}
\bibliography{references}

\end{document}
